% =====================================================================
\lezione{05/10/2026}{DPIA e data breach; presentazione del caso Arvenza}
  {Prof.~Giuseppe Vaciago}
  {regole del laboratorio; compliance anticipatoria; DPIA (contenuto, rischio inerente e residuo,
   consultazione preventiva); privacy by design e by default; data breach: definizione, artt.~33 e 34,
   consapevolezza del titolare, fattori di rischio, contenuto della notifica, responsabile del trattamento,
   remediation plan, registro dei data breach; presentazione del Lab~1 (Arvenza)}
\label{lez:2026-10-05}
% =====================================================================

\begin{nota}[Fonte]
Trascrizione automatica della registrazione, revisionata (pausa tra il minuto 94 e il 112 circa, con
formazione dei gruppi: non riportata). Slide di riferimento: \emph{02.~Cybersecurity Laws and
Regulations -- II Privacy} (moduli M4 \emph{DPIA} e M5 \emph{Data Breach}) e
\emph{03.~Data breach severity methodology~1.0} (ENISA, 2013). Il caso del laboratorio è nel PDF
\emph{Arvenza Breach Case Study} (vedi Capitolo~\ref{lab:01}).
\end{nota}

\section{Sintesi}
Lezione in due parti. La prima, più breve, è sulla \textbf{valutazione d'impatto (DPIA)}: il
docente la inquadra nel passaggio da una valutazione del rischio ``statica'' (basata sugli incidenti
già accaduti) a una \textbf{compliance anticipatoria}, che deve prevedere l'evoluzione del rischio
digitale. La seconda, centrale, è sul \textbf{data breach}: definizione, notifica al Garante entro 72
ore, comunicazione agli interessati, chi è ``il titolare'' in concreto, perché servono un piano prima
e un remediation plan dopo, e perché il \textbf{registro dei data breach} è la prova
dell'\emph{accountability}. In chiusura viene presentato il caso Arvenza per il laboratorio di
venerdì~9/10.

% ---------------------------------------------------------------------
\section{Organizzazione}
\begin{itemize}
  \item \textbf{Laboratori non obbligatori}, ma fanno parte del programma d'esame: durante il lab il
    docente dà indicazioni utili per le \textbf{domande aperte}. Chi non partecipa deve comunque
    ascoltare la registrazione (anche in differita; il docente suggerisce di usare un'IA per trovare
    i punti in cui parla, dato che lunghi tratti sono muti).
  \item \textbf{Punti bonus}: ogni laboratorio consegnato con esito sufficiente vale un piccolo
    punteggio aggiuntivo (in aula si parla di 0,25 punti a lab, su un totale previsto di 8--10 lab).
    Una parte dei gruppi (indicativamente il 25--30\%, forse di più) potrà presentare il proprio
    lavoro il lunedì successivo.
  \item \textbf{Formazione dei gruppi}: si è dedicata parte della lezione a creare i gruppi (in aula
    e online) per partire subito venerdì. Si può lavorare anche da soli, ma il docente lo sconsiglia.
  \item \textbf{Tempi del lab}: un'ora e mezza in aula, consegna via email al docente entro la fine
    del lab (indicate le 13:15); nessun lavoro richiesto prima o dopo. Solo per questo primo lab è
    consentito prepararsi in anticipo. Le consegne in ritardo non sono considerate, salvo situazioni
    eccezionali.
  \item \textbf{Formato libero}: il docente non fornisce un modello di severity assessment; ogni
    gruppo sceglie il formato che ritiene più efficace per presentare il rischio al Garante.
  \item Se pochi gruppi parteciperanno, il docente si riserva di cambiare formato ai laboratori.
  \item Il docente è disponibile per \textbf{tesi} (meno per tirocini, perché come avvocato ha meno
    contatti con aziende che cercano profili IT).
\end{itemize}

% ---------------------------------------------------------------------
\section{Dalla valutazione statica alla compliance anticipatoria}

Il settore in cui la valutazione del rischio è nata è la \textbf{salute e sicurezza sul lavoro}:
nell'industria (anche a Torino, nell'automotive) per decenni si moriva sul lavoro per mancanza di
prevenzione. In Italia il percorso arriva a compimento con il \textbf{D.Lgs.~81/2008}. In quel
contesto la valutazione può essere relativamente \emph{statica}: le macchine fanno sempre la stessa
cosa, l'errore nasce soprattutto dall'interazione uomo--macchina e la prevenzione passa in gran parte
dalla formazione.

Il ``corpo digitale'' è diverso: è ovunque e cambia a una velocità enorme (si pensi all'IA
generativa dal 2022 a oggi). Ne derivano tre conseguenze:
\begin{itemize}
  \item non si può risolvere il problema in modo \textbf{tattico} (``mettiamo l'antivirus''): serve una
    \textbf{strategia}, anche perché una PMI non ha il budget per prevenire tutto;
  \item il punto di rottura resta il \textbf{comportamento umano}, che va anticipato;
  \item la valutazione va scritta pensando a come verrà letta tra uno, due, tre anni:
    \textbf{compliance anticipatoria}, basata su elementi di \emph{foresight}. L'IA è insieme
    soluzione e fonte di nuovi rischi.
\end{itemize}

\begin{nota}[Indicazione per i laboratori]
Ogni valutazione fatta ``oggi'' deve reggere quando sarà riletta in futuro: chiedersi sempre come
apparirà quel documento tra qualche anno.
\end{nota}

% ---------------------------------------------------------------------
\section{La valutazione d'impatto (DPIA)}

\subsection{L'approccio del GDPR: fiducia nel titolare}
Il docente illustra la logica dell'\emph{accountability} con l'esempio di uno studente
(``Simone''): il GDPR si fida che il titolare faccia i compiti, senza che nessuno li controlli
preventivamente. La DPIA la redige il titolare e, se il rischio risulta alto, ha due strade: chiedere
la consultazione preventiva all'autorità o rinunciare al trattamento. Nella prassi, osserva il
docente, la tentazione è valutare il rischio come ``medio'' per evitare entrambe. Nell'AI Act invece,
secondo il docente, il sistema cambia: per i sistemi ad alto rischio non basta l'autovalutazione e
serve un passaggio con le autorità prima dell'uso (tema rinviato alla parte sull'AI Act).

\subsection{Quando serve una DPIA}
La DPIA va fatta quando un trattamento è \textbf{probabilmente ad alto rischio}, ma stabilirlo è
spesso una \textbf{zona grigia} lasciata alla valutazione del titolare (a differenza di norme come
NIS2, dove si entra o non si entra in una categoria). Esempi citati:
\begin{itemize}
  \item \textbf{videosorveglianza (CCTV)}: riconosciuta come trattamento che incide sui diritti
    fondamentali, quindi soggetta a DPIA;
  \item \textbf{whistleblowing}: chi lavora con un ente (dipendenti, collaboratori, fornitori) può
    segnalare violazioni in forma anonima; le piattaforme di whistleblowing sono soggette a DPIA.
\end{itemize}
Di conseguenza può capitare che per lo stesso strumento un'azienda faccia la DPIA e un'altra no,
perché i rispettivi DPO valutano il rischio in modo diverso.

\subsection{Contenuto e ciclo}
Dalle slide, la DPIA contiene: descrizione sistematica del trattamento e delle finalità;
valutazione di \textbf{necessità e proporzionalità}; valutazione dei rischi per i diritti e le libertà
degli interessati; misure previste per affrontarli e per dimostrare la conformità. Seguono
documentazione e monitoraggio.

Il rischio va valutato da \textbf{due punti di vista}: \textbf{giuridico} (impatto sui diritti
fondamentali) e \textbf{tecnico} (sicurezza del trattamento).

\paragraph{Necessità e proporzionalità.} Esempio del docente: chiedere a tutti gli studenti l'ultimo
certificato medico per dividerli in classi in base alle patologie può avere una logica sanitaria, ma
non è proporzionato, perché dà accesso a una quantità di informazioni sensibili che potrebbero essere
usate contro di loro.

\paragraph{Rischio inerente e residuo.}
\[
\text{rischio inerente} = \text{severità} \times \text{probabilità}
\qquad
\text{rischio residuo} = \text{rischio inerente} - \text{mitigazione}
\]
\begin{itemize}
  \item residuo \textbf{alto}: consultazione preventiva (art.~36) oppure stop del trattamento;
  \item residuo \textbf{medio o basso}: si prosegue, ma il docente raccomanda di ripetere la
    valutazione \textbf{ogni sei mesi}, perché nel digitale il rischio cambia in fretta (problema
    tipico: DPIA scritte nel 2020 e mai riviste).
\end{itemize}

\begin{esame}
Non fermarsi a ``rischio = severità $\times$ probabilità'': quello è solo il \textbf{rischio
inerente}. La valutazione si chiude sul \textbf{rischio residuo} (inerente meno mitigazioni), ed è il
residuo che decide se proseguire, consultare l'autorità o fermarsi. Il docente ha detto che lo
ripeterà anche in altri contesti.
\end{esame}

\paragraph{Natura e approccio (slide).} Valutazione \emph{a domande} (checklist, che il docente
preferisce perché evita dimenticanze), basata sul contesto; team interdisciplinare; analisi fin dalla
fase di progettazione; verifiche periodiche; approccio by design / by default. Una sola DPIA può
coprire trattamenti simili per natura, ambito, contesto, finalità e rischi. La pubblicazione non è
obbligatoria, ma l'EDPB suggerisce di pubblicarne almeno una sintesi o le conclusioni. Esiti possibili
della consultazione preventiva: ulteriori misure oppure stop del trattamento per mancanza di misure
di mitigazione disponibili.

\subsection{Privacy by design e by default (art.~25)}
\begin{itemize}
  \item \textbf{By design}: fin dall'inizio parte legale e parte tecnica lavorano insieme per
    integrare le garanzie nel trattamento, tenendo conto di stato dell'arte, costi di attuazione,
    natura, ambito, contesto, finalità e rischi (gli stessi criteri dell'\emph{accountability}).
  \item \textbf{By default}: per impostazione predefinita si trattano solo i dati necessari alla
    finalità (quantità, portata, conservazione, accessibilità). Contrasto con la cultura
    statunitense dell'\emph{opt-out}: in Europa vale l'\emph{opt-in}.
\end{itemize}

\begin{esame}
La DPIA non è, secondo il docente, una parte consistente dell'esame e una domanda specifica è
``rara'', ma \textbf{non si può escludere}; molte nozioni (proporzionalità, rischio residuo, by
design/by default) servono comunque per altre domande. Il ciclo della DPIA non va imparato a memoria.
\end{esame}

% ---------------------------------------------------------------------
\section{Data breach}

\begin{esame}
Sul data breach ci sarà \textbf{quasi sempre almeno una domanda} all'esame: è la parte più rilevante,
anche perché tutte le norme successive (DORA, AI Act, CRA, NIS2) prevedono notifiche di incidenti
costruite sullo stesso principio.
\end{esame}

\subsection{Perché il GDPR è stato il primo}
Il GDPR (deciso nel 2016, applicabile dal 2018) è stato la prima norma a imporre la notifica a
\textbf{tutti} i titolari, senza distinzione di settore: anche il pub che conserva le preferenze
alimentari dei clienti. Le norme successive sono invece segmentate (DORA solo per il settore
finanziario, NIS2 per alcune categorie). Anche dati apparentemente banali ``possono aprire molte
porte''.

\subsection{Definizione e tipologie}
L'\textbf{art.~4(12)} definisce la violazione dei dati personali come la violazione di sicurezza che
comporta accidentalmente o in modo illecito la \textbf{distruzione, perdita, modifica, divulgazione
non autorizzata o accesso} a dati personali trasmessi, conservati o comunque trattati.
Collegamento con la triade CIA:
\begin{itemize}
  \item distruzione e perdita $\rightarrow$ \textbf{disponibilità};
  \item modifica $\rightarrow$ \textbf{integrità};
  \item divulgazione o accesso non autorizzati $\rightarrow$ \textbf{riservatezza}.
\end{itemize}
Esempi dalle slide: smarrimento di un dispositivo non cifrato; ransomware (indisponibilità, ma spesso
anche riservatezza).

\begin{esame}
La definizione dell'art.~4(12) va \textbf{memorizzata bene}, insieme al collegamento con riservatezza,
integrità e disponibilità.
\end{esame}

\subsection{Gli obblighi: artt.~33 e 34}
\begin{itemize}
  \item \textbf{Art.~33}: notifica all'autorità di controllo \textbf{senza ingiustificato ritardo e,
    ove possibile, entro 72 ore}. Il ``ove possibile'' conta: se era materialmente impossibile (esempio
    del docente: incendio presso un fornitore) il ritardo può essere giustificato; se era possibile e
    non è stato fatto, è un problema.
  \item \textbf{Art.~34}: comunicazione agli \textbf{interessati} quando la violazione è suscettibile
    di presentare un \textbf{rischio elevato} per i loro diritti e libertà. Le aziende la temono per
    il danno reputazionale.
\end{itemize}

\subsection{Chi è ``il titolare'' quando parte l'orologio}
Le 72 ore decorrono da quando il \textbf{titolare è venuto a conoscenza} della violazione. Il titolare
è la persona giuridica, ma un'azienda ``è consapevole'' quando lo è chi la rappresenta: il
\textbf{legale rappresentante}.
\begin{itemize}
  \item Il legale rappresentante è nominato dai soci; nel 90\% dei casi coincide con il CEO, ma va
    verificato (in Italia con la \textbf{visura camerale}). Nelle multinazionali possono esistere
    deleghe specifiche, ad esempio in materia di trattamento dei dati.
  \item Per qualsiasi comunicazione legale a un'azienda conviene rivolgersi al legale
    rappresentante, che ha il dovere di rispondere.
  \item \textbf{Verità legale e verità reale}: se il legale rappresentante è informato per telefono
    martedì ma l'unica prova è un'email di giovedì, in giudizio conta ciò che si può provare.
    Da qui l'importanza delle \textbf{evidenze} della data di consapevolezza.
  \item Se il legale rappresentante non è raggiungibile (ferie, volo) e non c'è una procedura di
    delega, l'autorità lo leggerà come \textbf{carenza organizzativa}: secondo il docente in quel caso
    la sanzione è molto probabile.
\end{itemize}

\begin{esame}
Distinguere \textbf{errore umano} e \textbf{carenza organizzativa}. L'errore del singolo (il click sul
link di phishing) è umano e può essere compreso; la mancanza di procedure (chi avvisa chi, chi
decide quando il CEO è in volo) è un difetto prevedibile dell'organizzazione ed è ciò che le autorità
sanzionano. Strumento utile: la \textbf{matrice RACI} (\emph{Responsible, Accountable, Consulted,
Informed}), che sarà ripresa nelle esercitazioni su procedure e policy. Distinzione da usare nei lab.
\end{esame}

\subsection{Fattori di rischio per i diritti e le libertà}
Dalle slide, sei elementi da considerare come questione di metodo:
\begin{enumerate}[label=\alph*)]
  \item tipo di violazione;
  \item natura, numero e grado di sensibilità dei dati;
  \item facilità di associare i dati a una persona (es.\ dati pseudonimizzati: milioni di record
    ridotti a codici);
  \item gravità delle conseguenze (dati utilizzabili subito contro l'interessato: frode, sostituzione
    di persona);
  \item numero di interessati coinvolti;
  \item caratteristiche del titolare (un ospedale non è una casa vacanze).
\end{enumerate}
Il metodo operativo è la \textbf{metodologia ENISA} di severity assessment (vedi
Capitolo~\ref{lab:01} e la scheda al Capitolo~\ref{cap:schede}). Tutto è relativo: un milione di
interessati è molto per l'Italia, poco per l'India. Non esistono soluzioni bianco/nero.

\subsection{Contenuto della notifica (art.~33(3))}
Natura della violazione (categorie e numero approssimativo di interessati e di registrazioni);
contatti del \textbf{DPO}; probabili conseguenze; misure adottate o proposte per porre rimedio e
attenuarne gli effetti.

\begin{nota}[Attenzione alla sigla DPA]
``DPA'' indica sia la \emph{Data Protection Authority} sia il \emph{Data Processing Agreement}
(contratto con il responsabile ex art.~28). Il docente invita a chiedere chiarimenti, anche all'esame,
se la sigla è ambigua.
\end{nota}

Oltre ai contenuti, il titolare deve \textbf{motivare le proprie decisioni}, in particolare quando
decide di non notificare, ritarda la notifica o non comunica la violazione agli interessati. La
motivazione parte dal severity assessment: è l'applicazione diretta dell'\emph{accountability}.

\subsection{Lo scopo della notifica e la capacità di reazione}
In origine l'idea era lo \textbf{scambio di informazioni in tempo reale} per proteggere altre
organizzazioni da attacchi simili (il docente cita un lavoro del 2013, precedente al GDPR, e il ruolo
di centralizzazione di ENISA). Oggi la giurisprudenza delle autorità europee chiede di
\textbf{dimostrare la capacità di reazione}:
\begin{enumerate}
  \item un \textbf{piano prima}: procedure già in vigore (vedi il caso del legale rappresentante
    irraggiungibile);
  \item un \textbf{remediation plan dopo}. Alcuni esperti temono che un piano di rimedio equivalga
    ad ammettere una colpa, ma la violazione c'è già stata: più il piano è solido, più si abbassa il
    rischio. In alcuni casi il docente, in base all'accountability, ha deciso di non notificare
    perché il remediation plan riduceva il rischio in modo convincente.
\end{enumerate}

\begin{nota}[Indicazione per il laboratorio]
Il remediation plan può incidere sulla valutazione del rischio: considerarlo sempre, non limitarsi a
calcolare il numero. La slide sul remediation plan è indicata come ``molto importante'' per il lab.
\end{nota}

\subsection{Mappatura severità -- obblighi ed esenzioni dell'art.~34(3)}
Mappatura indicata in aula (e ripresa nel caso Arvenza):
\begin{center}
\begin{tabular}{lll}
\toprule
\textbf{Severità} & \textbf{Esito} & \textbf{Base} \\ \midrule
Bassa ($SE<2$) & solo registro & art.~33(5) \\
Media ($2 \le SE < 3$) & notifica all'autorità & art.~33 \\
Alta / molto alta ($SE \ge 3$) & autorità e interessati & art.~34, salvo 34(3) \\
\bottomrule
\end{tabular}
\end{center}
Il docente ammette che, da avvocato, si tende a ``tirare'' un po' questi livelli, ma vale ciò che è
scritto nella metodologia.

Esenzioni dalla comunicazione agli interessati (art.~34(3)):
\begin{enumerate}[label=\alph*)]
  \item misure tecniche e organizzative adeguate applicate ai dati violati (es.\ cifratura);
  \item misure successive che scongiurano il rischio elevato: difficile da sostenere, perché non si
    può garantire che una carta d'identità o di credito sottratta non venga usata;
  \item sforzo sproporzionato: si sostituisce con una \textbf{comunicazione pubblica}. Esempio vissuto
    dal docente: un'azienda di trasporto pubblico milanese con molti utenti senza email e non più
    rintracciabili, per cui si è scelto un avviso sui giornali.
\end{enumerate}

\subsection{Gestione del processo e costi}
\begin{itemize}
  \item \textbf{Preparazione}: piano di risposta agli incidenti, ruoli e responsabilità (incluso il
    CEO), simulazioni, modalità di interazione con le autorità. \textbf{Prevenzione}: formazione e
    misure di sicurezza.
  \item \textbf{Costi} di un data breach (slide): danno diretto, abbandono anomalo dei clienti, danno
    reputazionale, attività di help desk, indagini interne, spese legali, ripristino dei sistemi. Per il
    docente il più importante oggi è il \textbf{danno reputazionale}.
  \item \textbf{Aspetti pratici}: verificare quali notifiche sono obbligatorie; \textbf{analizzare con
    cura le ipotesi}, perché ogni assunzione può essere letta in due modi; usare i modelli e i canali
    ufficiali. Il modulo di notifica conta circa 55--60 domande: si è annunciata una simulazione di
    compilazione.
\end{itemize}

\subsection{Il responsabile del trattamento}
Il fornitore è oggi ``la base di tutti i problemi''. L'art.~28(3) rimanda al contratto; l'art.~33(2)
obbliga il responsabile a informare il titolare \textbf{senza ingiustificato ritardo} (nei contratti
si trovano tipicamente 24--36 ore). Nella pratica, soprattutto con i grandi fornitori, la
comunicazione arriva anche un mese dopo.

\begin{esame}
Se il responsabile avvisa con un mese di ritardo, il titolare non viola le 72 ore
\textbf{purché dimostri} che la violazione era imprevedibile e a lui ignota: l'orologio parte dalla
comunicazione del responsabile. Il titolare può poi chiedere i danni al responsabile per
inadempimento contrattuale. Tema ripreso nel Lab~2 (Capitolo~\ref{lab:02}).
\end{esame}

\subsection{Remediation plan (slide)}
Correggere la vulnerabilità che ha permesso l'incidente; \textbf{lesson learned}; condividere
informazioni con gli stakeholder per rafforzare il perimetro; se la causa è l'errore umano, rafforzare
la formazione. Sono esempi, non un elenco completo.

\subsection{Il registro dei data breach (art.~33(5))}
Strumenti: \textbf{incident log} (sintetico ma completo e aggiornato) e \textbf{report interni}
(modelli standard, decisioni documentate). L'art.~33(5) impone di documentare ogni violazione, le
circostanze, le conseguenze e i provvedimenti adottati, così che l'autorità possa verificare il
rispetto dell'articolo.

\begin{esame}
\textbf{Il registro è la prova dell'accountability.} Il controllo dell'autorità avviene su due livelli:
\begin{enumerate}
  \item \textbf{esistenza del registro}: se manca, la sanzione è automatica (e, se l'autorità trova i
    dati sul dark web, la massima possibile);
  \item \textbf{merito della valutazione}: se il rischio è stato dichiarato basso ma era evidentemente
    alto, si viene sanzionati comunque; la buona fede aiuta, ``me l'ha chiesto il CEO'' no.
\end{enumerate}
Il compito tipico dell'esperto di sicurezza sarà ridurre il rischio in modo che il titolare possa
legittimamente non notificare, ma sempre lasciando traccia del ragionamento.
\end{esame}

\paragraph{Firmare digitalmente il registro?} Su domanda di uno studente: un documento firmato
digitalmente (o con marca temporale, o inviato via PEC, o accompagnato da un hash) non può essere
modificato senza lasciare traccia e ``ferma l'orologio''. In Italia solo una minoranza di aziende
(5--10\% secondo il docente) lo fa; le altre tengono file Excel o PDF modificabili. Esistono anche
strumenti dedicati (es.\ OneTrust).

\paragraph{Alterare il registro è reato.} Le false dichiarazioni al Garante sono punite
dall'\textbf{art.~168 del Codice privacy} (reclusione da sei mesi a tre anni). Ne risponde non solo
il legale rappresentante ma anche chi ha materialmente alterato il documento: obbedire a un ordine
non scusa (\emph{ignorantia legis non excusat}). Se si deve correggere un file, meglio dirlo
all'autorità spiegando il motivo. Il docente racconta di un cliente che, dopo un'ispezione, gli
chiese di modificare il registro: rifiutò.

% ---------------------------------------------------------------------
\section{Presentazione del Lab~1: il caso Arvenza}
Il docente presenta il caso (dettagli e soluzione nel Capitolo~\ref{lab:01}): mailbox paghe
compromessa via phishing, sei gruppi di dati, tre posizioni in sala (CEO, IT Manager, HR Director).
Gli studenti fanno il \textbf{DPO}, che il docente descrive come funzione insieme legale e IT
(``nel mio sogno''), anche se in pratica è quasi sempre affidata a giuristi.
\begin{itemize}
  \item Formula ENISA: $SE = DPC \times EI + CB$.
  \item Va prodotto un \textbf{severity assessment} con un \textbf{parere} che concili o respinga le
    tre posizioni: ``non c'è una sola soluzione''.
  \item Non basta il numero: la mitigazione conta, e bisogna proporre una strategia (anche di
    prevenzione futura).
  \item Uso dell'IA: ammesso, ma se il lavoro è solo dell'IA sarà difficile presentarlo a voce.
\end{itemize}

% ---------------------------------------------------------------------
\section{Domande aperte e punti da verificare}
\begin{verifica}
\begin{itemize}
  \item \textbf{AI Act e autorizzazione preventiva}: la trascrizione fa dire al docente che per i
    sistemi ad alto rischio serve un'autorizzazione dell'autorità nazionale prima dell'uso. Il passaggio
    è confuso (si parla anche di FRIA) e va verificato sulle slide dell'AI Act.
  \item \textbf{Revisione semestrale della DPIA}: è una raccomandazione del docente, non un termine
    fissato dal GDPR (che chiede un riesame quando cambia il rischio, art.~35(11)). Confermare.
  \item Nella trascrizione il docente cita l'``art.~32.2'' per l'obbligo del responsabile di informare
    il titolare; le slide (e il testo del GDPR) indicano correttamente l'\textbf{art.~33(2)}.
  \item \textbf{Punti bonus}: il valore (0,25 a lab) e il numero di lab sono ricostruiti dal parlato,
    poco nitido; il totale massimo citato non è chiaro.
  \item Il ``paper del 2013'' sullo scambio di informazioni citato dal docente: autore e titolo non
    ricavabili con certezza dalla trascrizione.
  \item Concorso nel reato: il docente parla di ``associazione criminale'' tra legale rappresentante,
    IT manager ed esecutore; in termini tecnici potrebbe trattarsi di concorso di persone nel reato.
    Da non usare all'esame senza verifica.
\end{itemize}
\end{verifica}
